\begin{helper}[Terminal PE whose support first appears at the terminal stage]
Let \(\mathcal E\) be a positive-rung maximum ETT on a finite loopless
multigraph, with displayed tree \(T\), certified series \(\mathcal T\), and
rung number \(n\).  Let \(a=\operatorname{Fin.last}(n)\).  Suppose the data
at \(a\) certify a PE with chosen support \(v\), let \(h:v\in V(T)\), and let
\(b\in\operatorname{Fin}(n+1)\) satisfy
\[
 \noderef{ETTFirstAppearanceIndex}(\mathcal E,v,h)=b+1,
 \qquad b+1=n.
\]
If \(\sigma\) is proper, leaves exactly \(e_0\) uncolored, and is
\(\noderef{StableColoring}\) on the terminal ladder tree relative to
\((D_a,\varphi_a)\), then \(\sigma\) is
\(\noderef{StableColoringBeforeVertex}\) before \(v\) in \(T\), relative to
\((D_b,\varphi_b)\).  Moreover, if \(b>0\), it is stable on the ladder tree
at slot \(b\) relative to that same pair.
\end{helper}
\begin{proof}
First prove stability before \(v\).  Put \(T^*=T(v)-v\), whose vertices are precisely those specified by
\noderef{TreeVerticesBefore}.  The terminal PE certificate supplied by
\noderef{TashkinovSeries}, its support and ordered-prefix clauses in
\noderef{TashkinovSeriesStep}, and
\noderef{TreeSequenceInitialSegment} show that \(v\) and every vertex of
\(T^*\) belong to the terminal ladder tree \(T_n\).

Repeated use of \noderef{HkuMissingStoredMonotone} gives
\[
 \overline\varphi_{n-1}(T_n)\cup D_{n-1}
 \subseteq \overline\varphi_n(T_n)\cup D_n.                    \tag{1}
\]
If a color is missing on \(T^*\), \noderef{TreeMissingColors} gives a
vertex witness, and the preceding containment of vertex sets makes it
missing on \(T_n\).  Likewise, the endpoint witness in
\noderef{TreeIncidentEdge} makes every edge incident with \(T^*\) incident
with \(T_n\).  Thus the given stability relation restricts to a relation
from \(\sigma\) to \(\varphi_n\) on \(T^*\), with control \(D_{n-1}\), by
(1).  By the composition principle \noderef{sc2}, it suffices to prove that
\(\varphi_n\) is stable from \(\varphi_{n-1}\) on \(T^*\) with the same
control.

Eliminate the terminal certificate through the PE constructor of
\noderef{TashkinovSeriesStep}.  It supplies the selected coloring
\(\pi_{n-1}\), the second coloring \(\pi'_{n-1}\), the colors
\(\gamma_n,\delta_n\), stability of \(\pi_{n-1}\) from
\(\varphi_{n-1}\) with control \(D_{n-1}\), stability of
\(\pi'_{n-1}\) from \(\pi_{n-1}\) with control
\(D_{n-1}\cup\{\delta_n\}\), and
\[
 \varphi_n=\pi'_{n-1}/K,
\]
where \(K\) is the \((\gamma_n,\delta_n)\)-component through \(v\).  Its
\noderef{TreeExit} witness supplies a path \(R\), and
\noderef{TreeExitPath} says that \(R\) runs from outside \(T_n\) to \(v\)
and has no other vertex of \(T_n\).

We check that \(R\) exhausts \(K\).  The rule
\noderef{ParallelExtensionColor} makes \(\gamma_n\) missing at \(v\) under
\(\pi_{n-1}\), and the vertex clause of stability preserves this under
\(\pi'_{n-1}\).  The certificate \noderef{MaximumDefectiveChoice} gives an
edge \(f_n\) incident with \(v\) and colored \(\delta_n\) under
\(\pi_{n-1}\); the enlarged-control edge clause fixes this value under
\(\pi'_{n-1}\).  Properness makes \(f_n\) the unique two-color incidence at
\(v\).  At the outside endpoint of \(R\), one switch color is missing and
the path edge supplies the other, so properness again gives total degree
one.  Every internal path node already has its two path neighbors in
\noderef{TwoColorIncidenceGraph}.  If the connected support
\noderef{TwoColorComponent} contained an extra node, a shortest route to
\(R\) would give either a third neighbor to an internal node, contrary to
\noderef{TwoColorIncidenceDegreeAtMostTwo}, or a second neighbor to an
endpoint of total degree one.  Therefore \(R=K\), and \(K\) meets \(T_n\)
only at \(v\).

No edge incident with \(T^*\) has its edge-node in \(K\): otherwise its
endpoint in \(T^*\) would be adjacent to that edge-node in
\noderef{TwoColorIncidenceGraph} and would belong to the same
\noderef{TwoColorComponent}, contradicting the preceding intersection
statement.  Consequently \noderef{TwoColorComponentSwap} fixes every edge
incident with \(T^*\), and hence fixes every missing-color set there.
Restrict the two certified stability relations to \(T^*\), discard the
unused extra control color \(\delta_n\), and compose their edge equalities
and vertexwise missing-set equalities.  This proves stability of
\(\varphi_n\) from \(\varphi_{n-1}\) on \(T^*\).  A final application of
\noderef{sc2} with the restricted given relation proves the required
before-support stability.

For the additional ladder-tree conclusion, assume \(b>0\).  The equality
\(b+1=n\) and minimality in \noderef{ETTFirstAppearanceIndex} say that
\(v\) is absent from the ladder tree at slot \(b\) but belongs to the next
complete stage.  The certified step there and
\noderef{TreeSequenceInitialSegment} identify the predecessor tree entry by
entry with an initial segment of that next stage.  Consequently every
vertex of the predecessor tree occurs strictly before \(v\) in the final
display: an occurrence at or after \(v\)'s position would, by the
vertex-injectivity field of \noderef{TreeSequence}, either put \(v\) in the
predecessor or violate the initial-segment order.  Hence the predecessor
vertex set is contained in \(V(T^*)\).  By
\noderef{TreeIncidentEdge}, every edge incident with the predecessor is
therefore incident with \(T^*\); a missing color transports by the same
vertex witness through \noderef{TreeMissingColors}.  Restricting the edge
clause and vertexwise missing-set clause of the before-support stability
just proved gives stability on the ladder tree at slot \(b\), relative to
\((D_b,\varphi_b)\).  When \(b=0\) the implication has no premise, so no
stability assertion about the auxiliary seed is made.
\end{proof}
